\section{Instalação}

\begin{frame}
  \frametitle{Instalação}
  Docker Desktop está disponível para \textbf{Mac, Windows e Linux}.

  \vspace{0.5em}
  No Linux, provê pacotes \texttt{.deb} e \texttt{.rpm} para:
  \begin{itemize}
    \item Ubuntu, Debian, Red Hat e Fedora
  \end{itemize}

  \vspace{0.5em}
  \begin{alertblock}{Atenção}
    Existe versão experimental para Arch, porém sem suporte oficial.
  \end{alertblock}
\end{frame}

\begin{frame}
  \frametitle{Requisitos do Sistema — Linux}
  \begin{itemize}
    \item Kernel 64-bit com suporte a virtualização pela CPU
    \item Suporte a KVM para virtualização
    \item QEMU 5.2 ou superior
    \item Sistema de inicialização \texttt{systemd}
    \item 4 GB de RAM
  \end{itemize}

  \vspace{0.5em}
  \begin{alertblock}{Observação}
    O Docker Desktop no Linux roda uma máquina virtual.
  \end{alertblock}
\end{frame}

\begin{frame}
  \frametitle{Instalação no Ubuntu}
  Baixar o pacote mais recente e execute:

  \vspace{0.5em}
  \texttt{sudo apt-get update}

  \vspace{0.3em}
  \texttt{sudo apt install ./docker-desktop-amd64.deb}
\end{frame}